\chapter{Relacja wymuszania i twierdzenie o wymuszaniu}
\label{ch:forcing-theorem}
\chaptermark{Twierdzenie o wymuszaniu}

\section{Język wymuszania i gęstość lokalna}
W Rozdziale~\ref{ch:forcing-foundations} zdefiniowaliśmy
nazwy, ich wartości i tranzytywny zbiór $M[G]$.
Nie wykazaliśmy jeszcze, że $M[G]$ spełnia ZFC.
Zrobimy to poprzez relację wymuszania, definiowaną
w modelu wyjściowym, a następnie porównaną
z prawdziwością zdań w rozszerzeniu.

\begin{uwaga}[Założenia]
\label{rem:forcing33-background}
Przyjmujemy, że $M$ jest przeliczalnym
tranzytywnym modelem zbiorowym ZFC,
a $(\mathbb P,\leq,1_{\mathbb P})\in M$
jest pojęciem wymuszania z największym warunkiem.
Nie zakładamy antysymetrii ani separatywności.
Zawsze gdy piszemy o nazwach $\sigma,\tau,\rho$,
są to nazwy $\mathbb P$ należące do $M$;
$G$ jest filtrem $M$-generycznym.
Przeliczalność $M$ jest zewnętrzna.
Nie wywodzimy istnienia takiego modelu
z samej $\operatorname{Con}(\mathrm{ZFC})$.
\end{uwaga}

\begin{definicja}[Gęstość poniżej warunku]
\label{def:forcing33-local-dense}
Niech $p\in\mathbb P$, $D\subseteq\mathbb P$.
Zbiór $D$ jest \emph{gęsty poniżej $p$},
gdy dla każdego $q\leq p$ istnieje
$r\in D$ z $r\leq q$.
\end{definicja}

\begin{lemat}[Lokalna generyczność]
\label{lem:forcing33-local-generic}
Jeśli $p\in G$ i $D\in M$ jest gęsty
poniżej $p$, to $D\cap G\neq\varnothing$.
\end{lemat}
\begin{proof}
Niech $I_p=\{q\in\mathbb P:q\perp p\}$.
Zbiory $I_p$ i $D\cup I_p$ należą do $M$.
Jeżeli $q\perp p$, to $q\in I_p$.
W przeciwnym przypadku istnieje
$s\leq q,p$, a gęstość $D$ poniżej $p$
dostarcza $r\in D$ z $r\leq s\leq q$.
Zatem $D\cup I_p$ jest gęsty w $\mathbb P$.
Generyczność daje $t\in G\cap(D\cup I_p)$.
Ponieważ $p\in G$, każdy warunek $t\in G$
jest zgodny z $p$, więc $t\notin I_p$.
Stąd $t\in G\cap D$.
\end{proof}

\begin{przyklad}[Gęstość lokalna nie jest globalna]
\label{ex:forcing33-local}
W $\operatorname{Fn}(\omega,2)$ ustalmy
$p=\{(0,0)\}$ oraz
$D=\{q:q\leq p,\ 1\in\operatorname{dom}q\}$.
Każde wzmocnienie $p$ można przedłużyć
do elementu $D$. Jednak warunek
$\{(0,1)\}$ jest niezgodny z każdym
elementem $D$, więc $D$ nie jest
gęsty w całym wymuszaniu.
\end{przyklad}

\section{Atomowa relacja wymuszania}
W języku z parametrami będącymi nazwami
formuły atomowe mają postać $\sigma\in\tau$
albo $\sigma=\tau$. Nie definiujemy
wymuszania przez kwantyfikację po
rozszerzeniach generycznych: potrzebna
jest relacja \emph{wewnętrznie definiowalna}
w modelu $M$.

\begin{definicja}[Atomowe klauzule wymuszania]
\label{def:forcing33-atomic}
Niech $\sigma,\tau\in M$ będą nazwami,
a $p\in\mathbb P$. Przez jednoczesną
rekurencję ustalamy:
\[
\begin{aligned}
p\Vdash\sigma\in\tau
&\iff\forall q\leq p\ \exists r\leq q\
 \exists(\pi,s)\in\tau\
 (r\leq s\land r\Vdash\sigma=\pi),\\
p\Vdash\sigma=\tau
&\iff
 \bigl[\forall(\pi,s)\in\sigma\
  \forall q\leq p,s\
  \exists r\leq q\ (r\Vdash\pi\in\tau)\bigr]\\
&\quad{}\land
 \bigl[\forall(\pi,s)\in\tau\
  \forall q\leq p,s\
  \exists r\leq q\ (r\Vdash\pi\in\sigma)\bigr].
\end{aligned}
\]
W skrócie $\forall q\leq p,s$ wymagamy
jednocześnie $q\leq p$ i $q\leq s$.
\end{definicja}

\begin{uwaga}[Dobrze ufundowana rekurencja]
\label{rem:forcing33-ranks}
Dla nazwy $\tau$ niech $\rho(\tau)$ będzie
najmniejszym $\alpha$ z
$\tau\in N_{\alpha+1}$; hierarchia $N_\alpha$
pochodzi z Definicji~\ref{def:forcing-names}.
Każda pierwsza współrzędna $\pi$
pary $(\pi,s)\in\tau$ ma rangę
$\rho(\pi)<\rho(\tau)$.
Do uszeregowania dwóch argumentów
stosujemy sumę naturalną liczb
porządkowych $\alpha\mathbin{\#}\beta$,
zdefiniowaną przez dodawanie współczynników
przy jednakowych potęgach $\omega$
w postaci normalnej Cantora.
W odróżnieniu od zwykłego dodawania
jest ona ściśle rosnąca względem
każdego argumentu. W obu klauzulach
definicji zastąpienie jednej z nazw
jej pierwszą współrzędną zmniejsza
$\rho(\sigma)\mathbin{\#}\rho(\tau)$.
To dowodzi poprawności rekurencji
pozaskończonej dla \emph{par nazw},
również gdy nazwy są nieskończone.
\end{uwaga}

\begin{lemat}[Monotoniczność atomowego wymuszania]
\label{lem:forcing33-atomic-monotone}
Niech $\theta$ oznacza
$\sigma\in\tau$ albo $\sigma=\tau$.
Jeżeli $p\Vdash\theta$ i $q\leq p$,
to $q\Vdash\theta$.
\end{lemat}
\begin{proof}
Indukujemy jednocześnie po obu
klauzulach Definicji~\ref{def:forcing33-atomic}.
Każde wzmocnienie warunku $q$ jest
także wzmocnieniem $p$.
Wymagane przez klauzule świadectwa
istnieją więc już z założenia
$p\Vdash\theta$. Rekurencja po
rangach zapewnia, że przywołane
relacje pomocnicze są określone.
\end{proof}

\begin{lemat}[Domknięcie atomowe na gęstość]
\label{lem:forcing33-atomic-dense}
Niech $\theta$ będzie atomową formułą
z parametrami nazwami z $M$,
a $p\in\mathbb P$. Jeśli zbiór
$\{q\leq p:q\Vdash\theta\}$ jest
gęsty poniżej $p$, to $p\Vdash\theta$.
\end{lemat}
\begin{proof}
W przypadku $\theta=\sigma\in\tau$
dla każdego $q\leq p$ wybieramy
$r\leq q$ wymuszające przynależność.
Pierwsza klauzula dla $r$ dostarcza
dalszego wzmocnienia i pary z $\tau$,
a więc spełniona jest pierwsza
klauzula dla $p$.
Gdy $\theta=\sigma=\tau$, ustalmy
$(\pi,s)\in\sigma$ i $q\leq p,s$.
Gęstość daje $t\leq q$ wymuszające
$\sigma=\tau$. Z klauzuli równości
dla $t$ i warunku $t\leq s$
otrzymujemy $r\leq t$ wymuszające
$\pi\in\tau$. Analogicznie otrzymujemy
drugą połowę klauzuli równości.
\end{proof}

\section{Klauzule logiczne i definiowalność}
Każdą formułę pierwszego rzędu języka
teorii mnogości przedstawiamy przy użyciu
$=$, $\in$, $\neg$, $\land$ oraz $\exists$.
Pozostałe spójniki i $\forall$
są skrótami składniowymi.

\begin{definicja}[Wymuszanie formuł złożonych]
\label{def:forcing33-formulas}
Dla ustalonych formuł $\varphi,\psi$
z parametrami nazwami z $M$ określamy:
\[
\begin{aligned}
p\Vdash\neg\varphi
&\iff \neg\exists q\leq p\ (q\Vdash\varphi),\\
p\Vdash(\varphi\land\psi)
&\iff (p\Vdash\varphi)\land(p\Vdash\psi),\\
p\Vdash\exists x\,\varphi(x)
&\iff \forall q\leq p\ \exists r\leq q\
 \exists\tau\in M\
 \bigl(\tau\text{ jest nazwą}\land
                 r\Vdash\varphi(\tau)\bigr).
\end{aligned}
\]
Ostatni kwantyfikator po nazwach
oznacza kwantyfikację \emph{wewnątrz $M$}
po definiowalnej klasie nazw,
nie po utworzonym zbiorze wszystkich nazw.
\end{definicja}

\begin{lemat}[Monotoniczność i regularność wymuszania]
\label{lem:forcing33-regular}
Dla dowolnej ustalonej formuły
$\varphi(\vec\tau)$ z nazwami z $M$
oraz $p,q\in\mathbb P$:
\begin{enumerate}
\item z $q\leq p$ i $p\Vdash\varphi(\vec\tau)$
wynika $q\Vdash\varphi(\vec\tau)$;
\item jeśli
$\{q\leq p:q\Vdash\varphi(\vec\tau)\}$
jest gęsty poniżej $p$, to
$p\Vdash\varphi(\vec\tau)$.
\end{enumerate}
\end{lemat}
\begin{proof}
Pierwszy punkt dowodzimy indukcyjnie
po budowie formuły, zaczynając od
Lematu~\ref{lem:forcing33-atomic-monotone}.
Dla negacji zmniejszamy zakres
kwantyfikacji nad warunkami,
dla koniunkcji stosujemy indukcję
do obu składników, a w przypadku
egzystencji posługujemy się tą samą
gęstością poniżej każdego wzmocnienia.

Drugi punkt dla formuł atomowych jest
Lematem~\ref{lem:forcing33-atomic-dense}.
Dla koniunkcji jej wymuszanie
na zbiorze gęstym oznacza wymuszanie
każdego składnika na zbiorze gęstym;
hipoteza indukcyjna kończy dowód.
Jeżeli gęsto wiele warunków poniżej
$p$ wymusza $\neg\psi$, a istnieje
$r\leq p$ z $r\Vdash\psi$,
to jakieś $s\leq r$ wymusza negację.
Monotoniczność daje równocześnie
$s\Vdash\psi$, co przeczy definicji
negacji. Zatem $p\Vdash\neg\psi$.
Dla egzystencji: jeśli gęsto
wiele $r\leq p$ wymusza
$\exists x\,\psi(x)$, to dla każdego
$q\leq p$ możemy wybrać takie
$r\leq q$, a następnie $s\leq r$
i nazwę $\tau$ z
$s\Vdash\psi(\tau)$. Jest to
klauzula egzystencjalna dla $p$.
\end{proof}

\begin{twierdzenie}[Lemat definiowalności]
\label{thm:forcing33-definability}
Niech $M$ będzie tranzytywnym
modelem ZFC, a
$(\mathbb P,\leq,1_{\mathbb P})\in M$
pojęciem wymuszania.
Dla każdej \emph{ustalonej}
formuły pierwszego rzędu
$\varphi(x_1,\ldots,x_n)$ istnieje
formuła języka teorii mnogości
$\operatorname{Force}_\varphi
(p,\tau_1,\ldots,\tau_n,\mathbb P)$,
która w $M$ wyraża relację
$p\Vdash\varphi(\tau_1,\ldots,\tau_n)$
dla $p\in\mathbb P$ i nazw
$\tau_1,\ldots,\tau_n\in M$.
W konsekwencji zbiory warunków
opisane za pomocą ustalonej relacji
wymuszania i parametrów z $M$
należą do $M$.
\end{twierdzenie}
\begin{proof}
Rekurencja atomowa
z Definicji~\ref{def:forcing33-atomic}
jest dobrze ufundowana dzięki
Uwadze~\ref{rem:forcing33-ranks}.
Dla ustalonej pary nazw można
zamknąć zbiór obu nazw na
występujące dziedzicznie pierwsze
współrzędne par. To jest zbiór
należący do $M$ (wykorzystujemy
rekurencję, sumę i zastępowanie).
Rekurencja po odpowiednich parach
nazw i wszystkich warunkach
$p\in\mathbb P$ tworzy w $M$
jednoznaczny wynik. Dla danej pary
zdanie „istnieje poprawny wynik
tej rekurencji, w którym wskazana
relacja zachodzi” jest formułą
pierwszego rzędu. Jednoznaczność
gwarantuje zgodność wyników
uzyskanych z różnych ograniczonych
zbiorów nazw. Tak otrzymujemy
definiowalność relacji atomowych.

Następnie przeprowadzamy indukcję
po \emph{ustalonej formule}
$\varphi$. Negacja i koniunkcja
używają jedynie znanych już
formuł definiujących wymuszanie
składników. W przypadku
$\exists x\,\psi(x)$ kwantyfikator
po nazwach jest kwantyfikatorem
pierwszego rzędu wewnątrz $M$
z definiowalnym warunkiem
bycia nazwą. Otrzymujemy zatem
formułę $\operatorname{Force}_\varphi$.
Wyróżnianie wewnątrz $M$
tworzy wymagane podzbiory
$\mathbb P$.
\end{proof}

\begin{uwaga}[Schemat, nie predykat prawdy]
\label{rem:forcing33-tarski}
Lemat definiowalności dotyczy
\emph{każdej ustalonej formuły}
oddzielnie; nie jest wewnętrzną
definicją jednego predykatu prawdy
dla wszystkich kodów zdań.
Dla wymuszania jednoelementowego
ten ostatni postulat przeczyłby
twierdzeniu Tarskiego o
niedefiniowalności prawdy.
\end{uwaga}

\section{Lemat prawdy i twierdzenie o wymuszaniu}

\begin{lemat}[Gęstość rozstrzygnięć]
\label{lem:forcing33-decision}
Dla ustalonej formuły $\varphi(\vec\tau)$
z nazwami z $M$ zbiór
\[
D_\varphi=\{p\in\mathbb P:
p\Vdash\varphi(\vec\tau)
\text{ lub }p\Vdash\neg\varphi(\vec\tau)\}
\]
jest gęsty w $\mathbb P$ i należy do $M$.
\end{lemat}
\begin{proof}
Przynależność do $M$ wynika z
Twierdzenia~\ref{thm:forcing33-definability}.
Dla danego $p\in\mathbb P$ albo
istnieje $q\leq p$ wymuszające
$\varphi$, albo żadne wzmocnienie
$p$ nie wymusza $\varphi$.
W pierwszym przypadku $q\in D_\varphi$,
w drugim $p\Vdash\neg\varphi$,
więc $p\in D_\varphi$.
\end{proof}

\begin{twierdzenie}[Atomowy lemat prawdy]
\label{thm:forcing33-atomic-truth}
Niech $\sigma,\tau\in M$ będą
nazwami, a $G$ filtrem $M$-generycznym.
Dla $\theta$ postaci $\sigma\in\tau$
lub $\sigma=\tau$ zachodzi
\[
M[G]\models\theta(\sigma^G,\tau^G)
\iff\exists p\in G\
(p\Vdash\theta(\sigma,\tau)).
\]
\end{twierdzenie}
\begin{proof}
Przeprowadzamy jednoczesną indukcję
po sumie naturalnej rang par nazw
z Uwagi~\ref{rem:forcing33-ranks}.

\emph{Przynależność, od prawej do lewej.}
Jeśli $p\in G$ wymusza $\sigma\in\tau$,
to zbiór $q\leq p$, dla których istnieje
$(\pi,s)\in\tau$ z
$q\leq s$ i $q\Vdash\sigma=\pi$,
jest gęsty poniżej $p$ i należy do $M$.
Z Lematu~\ref{lem:forcing33-local-generic}
wybieramy taki $q\in G$.
Ponieważ $q\leq s$, mamy $s\in G$.
Indukcja dla równości daje
$\sigma^G=\pi^G\in\tau^G$.

\emph{Przynależność, od lewej do prawej.}
Jeśli $\sigma^G\in\tau^G$, to
istnieje $(\pi,s)\in\tau$
z $s\in G$ i $\pi^G=\sigma^G$.
Indukcja dla równości daje
$p\in G$ wymuszające $\sigma=\pi$.
Wspólne wzmocnienie $q\in G$
warunków $p,s$ wymusza $\sigma=\pi$.
Każde $r\leq q$ pozostaje
silniejsze od $s$ i wymusza
tę samą równość; zatem
$q\Vdash\sigma\in\tau$.

\emph{Równość, od prawej do lewej.}
Jeśli $p\in G$ wymusza
$\sigma=\tau$ i
$x=\pi^G\in\sigma^G$,
wybieramy $(\pi,s)\in\sigma$,
$s\in G$, i wspólne
wzmocnienie $q\in G$ poniżej $p,s$.
Klauzula równości daje gęstość
poniżej $q$ zbioru warunków
wymuszających $\pi\in\tau$.
Z lokalnej generyczności
otrzymujemy taki warunek w $G$;
indukcja daje $x\in\tau^G$.
Symetrycznie $\tau^G\subseteq\sigma^G$.

\emph{Równość, od lewej do prawej.}
Załóżmy $\sigma^G=\tau^G$.
Definiujemy $E\in M$ jako zbiór
warunków $p$ spełniających jeden
z warunków:
\begin{enumerate}
\item $p\Vdash\sigma=\tau$;
\item istnieje $(\pi,s)\in\sigma$
z $p\leq s$ i bez żadnego
$q\leq p$ wymuszającego $\pi\in\tau$;
\item istnieje $(\pi,s)\in\tau$
z $p\leq s$ i bez żadnego
$q\leq p$ wymuszającego $\pi\in\sigma$.
\end{enumerate}
Jest to zbiór definiowalny,
a jego gęstość wynika z definicji:
gdy $p$ nie wymusza równości,
zanegowanie jednej z dwóch klauzul
daje jej świadectwo i wzmocnienie
$q\leq p$ należące do drugiego
lub trzeciego członu $E$.
Wybieramy $p\in G\cap E$.
W przypadku drugim z $p\leq s$
wynika $s\in G$, zatem
$\pi^G\in\sigma^G=\tau^G$.
Indukcja dla przynależności
daje $q\in G$ wymuszające
$\pi\in\tau$. Wspólne wzmocnienie
$p,q$ przeczy drugiemu warunkowi
definiującemu $E$. Trzeci
przypadek odrzucamy analogicznie.
Pozostaje $p\Vdash\sigma=\tau$.
\end{proof}

\begin{twierdzenie}[Lemat prawdy]
\label{thm:forcing33-truth}
Dla ustalonej formuły
$\varphi(x_1,\ldots,x_n)$, nazw
$\tau_1,\ldots,\tau_n\in M$
oraz filtru $M$-generycznego $G$
zachodzi równoważność
\[
M[G]\models\varphi(\tau_1^G,\ldots,\tau_n^G)
\iff\exists p\in G\
(p\Vdash\varphi(\tau_1,\ldots,\tau_n)).
\]
\end{twierdzenie}
\begin{proof}
Indukujemy po budowie formuły.
Podstawą jest
Twierdzenie~\ref{thm:forcing33-atomic-truth}.
Dla koniunkcji z prawdziwości
obu składników otrzymujemy warunki
$p,q\in G$ wymuszające składniki,
a ich wspólne wzmocnienie wymusza
koniunkcję. Implikacja odwrotna
jest bezpośrednia.

Gdy $p\in G$ wymusza $\neg\psi$,
prawdziwość $\psi$ w $M[G]$
prowadziłaby do warunku $q\in G$
wymuszającego $\psi$.
Wspólne wzmocnienie $p,q$
przeczy klauzuli negacji.
Odwrotnie, jeśli $M[G]\models\neg\psi$,
generyczność daje $p\in G\cap D_\psi$
z Lematu~\ref{lem:forcing33-decision}.
Nie może on wymuszać $\psi$,
więc wymusza $\neg\psi$.

Jeśli $M[G]\models\exists x\,\psi(x,\vec\tau)$,
świadek ma postać $\rho^G$ dla nazwy
$\rho\in M$. Indukcja daje
$p\in G$ wymuszające $\psi(\rho,\vec\tau)$.
Każde wzmocnienie $p$ wymusza tę
formułę, więc $p$ wymusza egzystencję.
Odwrotnie, gdy $p\in G$ wymusza
$\exists x\,\psi(x,\vec\tau)$, zbiór
\[
D=\{q\leq p:\exists\rho\text{ nazwa w }M\
(q\Vdash\psi(\rho,\vec\tau))\}
\]
należy do $M$ i jest gęsty poniżej $p$.
Lemat~\ref{lem:forcing33-local-generic}
daje $q\in G\cap D$ oraz nazwę $\rho$
wymuszającą odpowiednią formułę.
Indukcja daje
$M[G]\models\psi(\rho^G,\vec\tau^G)$,
a więc prawdziwość egzystencji.
\end{proof}

\begin{twierdzenie}[Pełne twierdzenie o wymuszaniu]
\label{thm:forcing33-main}
Niech $M$ będzie przeliczalnym tranzytywnym
modelem ZFC i
$(\mathbb P,\leq,1_{\mathbb P})\in M$
pojęciem wymuszania. Dla każdej
ustalonej formuły $\varphi(\vec x)$,
nazw $\vec\tau\in M$ oraz $p\in\mathbb P$:
\begin{enumerate}
\item relacja $p\Vdash\varphi(\vec\tau)$
jest definiowalna w $M$;
\item dla każdego filtru $M$-generycznego $G$
zachodzi lemat prawdy:
$M[G]\models\varphi(\vec\tau^G)$
dokładnie wtedy, gdy pewien $q\in G$
wymusza $\varphi(\vec\tau)$;
\item zachodzi równoważność semantyczna
\[
 p\Vdash\varphi(\vec\tau)
 \iff \forall G\,
 \bigl(G\text{ jest $M$-generyczny},\
 p\in G\ \Longrightarrow\
 M[G]\models\varphi(\vec\tau^G)\bigr).
\]
\end{enumerate}
\end{twierdzenie}
\begin{proof}
Punkty pierwszy i drugi to
Twierdzenia~\ref{thm:forcing33-definability}
i~\ref{thm:forcing33-truth}.
Implikacja w prawo w punkcie trzecim
wynika z lematu prawdy.
Jeżeli $p$ nie wymusza formuły,
to na mocy Lematu~\ref{lem:forcing33-regular}
zbiór warunków ją wymuszających
nie jest gęsty poniżej $p$.
Istnieje więc $q\leq p$, poniżej
którego nikt nie wymusza tej formuły,
czyli $q\Vdash\neg\varphi(\vec\tau)$.
Wniosek~\ref{cor:countable-generic}
zapewnia filtr $M$-generyczny $G$
zawierający $q$, a więc i $p$.
Z lematu prawdy wynika
$M[G]\models\neg\varphi(\vec\tau^G)$.
To przeczy prawej stronie równoważności.
\end{proof}

\begin{uwaga}[Brak błędnego koła]
\label{rem:forcing33-no-circularity}
W dowodzie lematu prawdy i twierdzenia
o wymuszaniu nie korzystaliśmy z
$M[G]\models\mathrm{ZFC}$.
Znaliśmy tylko tranzytywny
zbiór $M[G]$ z Rozdziału~\ref{ch:forcing-foundations}
i zwykłą semantykę pierwszego rzędu
dla struktury $(M[G],\in)$.
Dopiero poniżej zastosujemy otrzymane
twierdzenie do sprawdzenia aksjomatów ZFC.
\end{uwaga}

\begin{przyklad}[Rozstrzygnięcie bitu Cohena]
\label{ex:forcing33-cohen-bit}
W $\operatorname{Fn}(\omega,2)$
niech $\dot C$ będzie nazwą
generycznego podzbioru $\omega$
z Przykładu~\ref{ex:forcing-cohen-name}.
Jeżeli $n\in\operatorname{dom}p$,
to $p$ wymusza $\check n\in\dot C$
wtedy i tylko wtedy, gdy $p(n)=1$;
gdy $p(n)=0$, wymusza negację.
Każdy filtr generyczny zawierający $p$
daje tę samą wartość $c_G(n)=p(n)$,
a więc wniosek jest konsekwencją
Twierdzenia~\ref{thm:forcing33-main}.
Jeżeli $n\notin\operatorname{dom}p$,
istnieją niezgodne wzmocnienia
$p\cup\{(n,0)\}$ i $p\cup\{(n,1)\}$.
Mają one przedłużenia do filtrów
generycznych, które nadają
różne wartości współrzędnej $n$;
sam $p$ nie rozstrzyga tej formuły.
\end{przyklad}

